\chapter*{Thesis Conclusion}
\addcontentsline{toc}{chapter}{Thesis Conclusion}


% \paragraph{Objectives revisited.}
This dissertation investigated how to satisfy information needs in knowledge-intensive legal and investigative contexts, considering use cases such as precedent retrieval, document navigation, investigative retrieval, question answering, and statistical analysis. The work pursued three objectives: (\ref{ob:extraction_assessment}) assessing the applicability and limitations of general-domain \acl{ee} methods in these domains; (\ref{ob:architecture}) designing an entity-centric \acl{dia} that enables traceability and remains useful despite imperfect extraction; and (\ref{ob:refactx}) enabling efficient and scalable \acl{qa} which can be executed on local hardware and whose answers can be verified by reviewing source information. These objectives are motivated by the need to access the information consolidated from heterogeneous, high-risk settings while retaining traceability and human oversight.

% \paragraph{Synthesis of contributions.}
\medskip\noindent
\Cref{chap:extraction} showed that incremental \acl{ee}---in which unknown entities, such as the persons mentioned in chat messages, are recognized and added to a \acl{kr} as novel entities---suffers from error propagation and the NIL prediction (\ie the detection of novel entity mentions) is a major source of errors. Moreover, while general-domain \acl{ee} models can be applied to legal documents and investigative chat data, in-domain fine-tuning or \acl{hitl} correction may be required for achieving satisfactory extraction quality. 

\Cref{chap:architecture} introduced an entity-centric architecture for integrating heterogeneous data (judgments, investigative chats, attachments) that remains operational under imperfect extraction by enabling users to trace automatic outputs to source information and to correct them in a \acl{hitl} setting. The design is grounded in real tasks and prior prototypes, and it provides the abstractions required to support retrieval, navigation, \acl{qa}, and statistical analysis over such data. In addition, the chapter presented advanced user interfaces that can be coupled with an implementation of the architecture and operationalize the use cases addressed in this thesis.

\Cref{chap:refactx} showed that efficient and scalable \acl{qa} with grounded answers can be achieved by using an \ac{llm} with constrained generation supported by a disk-backed prefix tree, without relying on retrievers, complex pipelines, or architectural changes to the \ac{llm}, which may be impractical under budget or deployment constraints. The chapter introduced \refactx, a constrained-generation approach that produces valid facts from large \acp{kb}, enabling users to verify that the answer is consistent with the acquired grounded facts. The method scales to hundreds of millions of facts with negligible latency and yields substantial gains over \acp{llm} relying solely on parametric knowledge.


Taken together, these results address the considered \acl{ia} use cases for satisfying the information needs of users in knowledge-intensive domains, with special focus on sensitive contexts, such as law and investigations, for which additional warranties are proposed---including traceability, verifiability, and error correction capability. In fact, the \acl{dia} presented in \Cref{chap:architecture} supports the integration of extraction services, such as the ones evaluated in \Cref{chap:extraction}, and---by using a \acl{kb} supporting entity facts---of \refactx for answering questions on extracted facts, such as the ones related to the suspect's communication network described in \Cref{sec:eechats}. 

% \paragraph{Outlook.} % future work?
\medskip
\noindent
Several directions remain open. First, extraction could be improved by advancing NIL prediction and clustering, by adapting \acl{kc} models with in-domain data, and by evaluating newer approaches and \acp{llm}, with particular interest in adaptation techniques, such as \acl{icl}~\cite{jm3}, which require minimal manual operations. Second, the architectural design could be implemented and evaluated, as well as advanced \acp{ui} for currently uncovered use cases such as statistical analysis (\ref{uc:stats}). Third, richer \acl{hitl} feedback loops could be explored, where user corrections are not only recorded but also incorporated to update models or constraints, thereby reducing recurrent errors. Finally, methods for extracting facts from unstructured text, while preserving traceability, could be explored to better integrate \refactx into the \acl{dia} in scenarios where data are predominantly unstructured.

\section*{Ethical and regulatory considerations}
\addcontentsline{toc}{section}{Ethical and regulatory considerations}
The design choices of this thesis align with the constraints imposed on the use of \ac{ai} in sensitive domains. In legal and investigative contexts, data typically contain personal or confidential information and cannot be transferred to third-party \ac{api}-based services without risking non-compliance with the \emph{GDPR (Regulation (EU) 2016/679)}~\cite{gdpr2016}, which requires lawful and controlled processing~\cite[Art.~5]{gdpr2016} and restricts international data transfers~\cite[Art.~44--49]{gdpr2016}. The \emph{AI Act (Regulation (EU) 2024/1689)}~\cite{ai_act} further classifies \ac{ai} used in law enforcement and justice as high-risk~\cite[Annex~III]{ai_act}, requiring risk management, data governance, logging, and traceability~\cite[Art.~9--12]{ai_act}.
Using \acp{llm} via external \acp{api} introduces compliance complexities; for instance, a data
processing agreement must be in place between controller and processor~\cite[Art.~28(3)]{gdpr2016}.
By contrast, the approaches developed in this thesis---notably the local execution of \acl{qa}
and the traceability, verifiability, and error correction ensured by the \acl{dia}---keep processing within the
secure perimeter of the deploying institution and preserving human oversight~\cite[Art.~14]{ai_act}.


